\chapter{Programming for FRAME21}
\label{ch:frame21}

\section{Introduction}
\topic{Purpose}
Develop \textsc{FRAME21} --- a linear analysis program for framed structures.

\topic{Procedure}
\begin{enumerate}
  \item Input data.
  \item Set up global degrees of freedom (equation numbering).
  \item Assemble external loads into the load vector.
  \item Form the global stiffness matrix and add the equivalent forces (arising from the fixed-end forces) to the load vector.
  \item Solve for the nodal unknowns.
  \item Calculate member displacements and member end forces (force recovery).
\end{enumerate}

\section{Input data}
Functions: \texttt{INPUT} and \texttt{OUTPUT1}; read in and write out the following data.

\begin{enumerate}[leftmargin=*]
  \item \textbf{Structural data}
    \begin{description}[font=\ttfamily,leftmargin=6em,style=sameline,itemsep=0pt]
      \item[TITLE] title of the problem analyzed; e.g.\ Example 5.3 of McGuire et al.\ (2000).
      \item[FUNIT] unit of force such as kN, kips.
      \item[LUNIT] unit of length such as mm, in.
      \item[NNOD] number of nodes
      \item[NBC] number of elements (number of beam columns)
      \item[NMAT] number of sets of material properties
      \item[NSEC] number of sets of sectional properties
      \item[ITP] frame type; 1: beam, 2: 2D truss, 3: 2D frame, 4: grid, 5: 3D truss, 6: 3D frame.
      \item[NCO] number of coordinates per node
      \item[NDN] number of d.o.f.\ per node
      \item[NNE] number of nodes per element
      \item[NDE] number of d.o.f.\ per element
      \item[IFORCE] indicator for inputting the fixed-end forces (FEF). When \texttt{IFORCE}=1, no FEF is needed; when \texttt{IFORCE}=2, FEF needs to be input.
    \end{description}

  \item \textbf{Nodal data}
    \begin{enumerate}[label=(\alph*)]
      \item Coordinates: \texttt{COOR(NCO,NNOD)} --- $x$, $y$, and $z$.
      \item Boundary conditions: \texttt{NFIX(NDN,NNOD)} --- $u$, $v$, $w$, $\theta_x$, $\theta_y$, and $\theta_z$ (3D frame).

        Specify ``$-1$'' if the d.o.f.\ is constrained (or fixed); specify ``0'' if free. Specify the smaller node number for the d.o.f.\ which is the same as the d.o.f.\ of the larger node number if the \emph{double node technique} is used. For example, if nodes 3 and 4 are double nodes (they have the same coordinates) used to model an internal pin for a beam, then the \texttt{NFIX} values are 0, 0 for node 3, and 3, 0 for node 4.
        \begin{fig}
          % Beam with an internal pin modelled by double nodes 3 and 4
\begin{tikzpicture}[scale=1.0,font=\small]
  \draw[member] (0,0) -- (7,0);
  \pic at (0,0) {wallL}; \pic at (7,0) {wallR};
  \pic at (2.3,0) {roller};
  \node[node dot] at (4.6,0) {};
  \node[above right,inner sep=1pt] at (0,0) {1};
  \node[above right,inner sep=1pt] at (2.3,0.02) {2};
  \node[above left,inner sep=1pt] at (4.5,0.05) {3};
  \node[above right,inner sep=1pt] at (4.7,0.05) {4};
  \node[above left,inner sep=1pt] at (7,0) {5};
  % rotations
  \foreach \x/\l in {2.3/2,4.35/3,4.95/4}{
    \draw[->,thick] ($(\x,0)+(0.3,0.35)$) arc (20:160:0.3);
    \node[above] at (\x,0.6) {$\theta_\l$};}
  \draw[->,thick] (4.6,-0.9) -- (4.6,-0.15) node[pos=0,right] {$v$};
\end{tikzpicture}

        \end{fig}
        \begin{center}
          \begin{minipage}{0.47\linewidth}\centering
            \texttt{NFIX(NDN,NNOD)}\\[4pt]
            \begin{tabular}{|c|c|c|c|c|c|}\hline
                       & 1    & 2    & 3 & 4 & 5   \\\hline
              $v$      & $-1$ & $-1$ & 0 & 3 & $-1$\\\hline
              $\theta$ & $-1$ & 0    & 0 & 0 & $-1$\\\hline
            \end{tabular}
          \end{minipage}
          \begin{minipage}{0.5\linewidth}\centering
            \texttt{IDND(NDN,NNOD)} (discussed later)\\[4pt]
            \begin{tabular}{|c|c|c|c|c|c|}\hline
                       & 1    & 2    & 3 & 4 & 5   \\\hline
              $v$      & $-1$ & $-3$ & 2 & 2 & $-4$\\\hline
              $\theta$ & $-2$ & 1    & 3 & 4 & $-5$\\\hline
            \end{tabular}
          \end{minipage}
        \end{center}
      \item External ``concentrated'' nodal loads: \texttt{EXLD(NDN,NNOD)}
        \begin{fig}
          % Portal frame with concentrated nodal loads
\begin{tikzpicture}[scale=1.0,font=\small]
  \draw[member] (0,0) -- (0,3) -- (3.2,3) -- (3.2,0);
  \pic at (0,0) {fixed}; \pic at (3.2,0) {pin};
  \node[right,inner sep=2pt] at (0,0.3) {1}; \node[below right,inner sep=2pt] at (0,3) {2};
  \node[below left,inner sep=2pt] at (3.2,3) {3}; \node[above left,inner sep=2pt] at (3.2,0.05) {4};
  \draw[load] (-1.0,3) -- (-0.1,3) node[pos=0,left] {2};
  \draw[load] (0,4.0) -- (0,3.1) node[pos=0,above] {10};
  \draw[load] (3.2,3.1) -- (3.2,4.0) node[above] {5};
  \draw[->,thick] ($(3.2,3)+(0.45,0.2)$) arc (30:-170:0.45); \node[right] at (3.75,2.95) {3};
  \draw[load] (4.2,0) -- (3.4,0) node[pos=0,right] {7};
\end{tikzpicture}

        \end{fig}
        \begin{center}
          \begin{minipage}{0.47\linewidth}\centering
            \texttt{NFIX(NDN,NNOD)}\\[4pt]
            \begin{tabular}{|c|c|c|c|c|}\hline
                       & 1    & 2 & 3 & 4   \\\hline
              $u$      & $-1$ & 0 & 0 & 0   \\\hline
              $v$      & $-1$ & 0 & 0 & $-1$\\\hline
              $\theta$ & $-1$ & 0 & 0 & 0   \\\hline
            \end{tabular}
          \end{minipage}
          \begin{minipage}{0.5\linewidth}\centering
            \texttt{EXLD(NDN,NNOD)}\\[4pt]
            \begin{tabular}{|c|c|c|c|c|}\hline
                    & 1 & 2     & 3    & 4   \\\hline
              $F_x$ & 0 & 2     & 0    & $-7$\\\hline
              $F_y$ & 0 & $-10$ & 5    & 0   \\\hline
              $M$   & 0 & 0     & $-3$ & 0   \\\hline
            \end{tabular}
          \end{minipage}
        \end{center}
    \end{enumerate}

  \item \textbf{Elemental data}
    \begin{enumerate}[label=(\alph*)]
      \item Identification matrices (connectivities):\\
        \texttt{IDBC(5,NBC)}: node 1, node 2, material type, sectional type, blank.
      \item Direction cosines of the local $y$ axis: \texttt{VECTY(3,NBC)}. (\texttt{ITP}=6 only)
        \begin{fig}
          % Global (x,y,z) and local (x',y',z') axes of a 3D frame member
\begin{tikzpicture}[scale=1.0,font=\small]
  \begin{scope}
    \draw[->,thick] (0,0) -- (0,2.0) node[above] {$z$};
    \draw[->,thick] (0,0) -- (-1.5,-0.9) node[left] {$x$};
    \draw[->,thick] (0,0) -- (1.5,-0.9) node[right] {$y$};
  \end{scope}
  \begin{scope}[xshift=6cm]
    % a small box (member cross-section sketch)
    \draw (0,0) -- (0.5,0) -- (0.5,1.3) -- (0,1.3) -- cycle;
    \draw (0,1.3) -- (-0.4,1.55) -- (0.1,1.55) -- (0.5,1.3);
    \draw (0,0) -- (-0.4,0.25) -- (-0.4,1.55);
    \draw (-0.4,1.55) -- (-1.4,2.1); \draw (0.1,1.55) -- (-0.9,2.1);
    \draw[->,thick] (0,0) -- (0,2.4) node[above] {$y'$};
    \draw[->,thick] (0,0) -- (-1.6,-0.7) node[left] {$z'$};
    \draw[->,thick] (0,0) -- (1.6,-0.7) node[right] {$x'$};
    \node[below,font=\footnotesize] at (0,-1.1) {$0,0,1$ for this example};
  \end{scope}
\end{tikzpicture}

        \end{fig}
      \item Fixed-end forces (in local axis) \texttt{FEF(NDE,NBC)} $= \F_l$. (\texttt{IFORCE}=2 only)
    \end{enumerate}

  \item \textbf{Material and sectional properties}
    \begin{enumerate}[label=(\alph*)]
      \item \texttt{PROP(5,NMAT)}: Young's modulus $E$, Poisson's ratio $\nu$, blank, blank, blank.\\
        For dynamic problems, $\rho$ is required; for plastic analysis problems, $\sigma_y$ is needed.
      \item \texttt{SECT(5,NSEC)}: cross-sectional area $A$, moments of inertia $I_z$ and $I_y$, and torsional constant $J$, blank.
    \end{enumerate}
\end{enumerate}

\subsection*{Example of input files}
\begin{fig}
  % 2-D truss example for the input file (8 nodes, 14 members)
\begin{tikzpicture}[scale=0.55,font=\small]
  \coordinate (n1) at (0,0); \coordinate (n2) at (6,0); \coordinate (n3) at (12,0); \coordinate (n4) at (18,0); \coordinate (n5) at (24,0);
  \coordinate (n6) at (6,8); \coordinate (n7) at (12,8); \coordinate (n8) at (18,8);
  \draw[thinmember] (n1) -- (n5); \draw[thinmember] (n6) -- (n8);
  \draw[thinmember] (n1) -- (n6) -- (n2); \draw[thinmember] (n6) -- (n3) -- (n7); \draw[thinmember] (n3) -- (n8) -- (n4);
  \draw[thinmember] (n4) -- (n7); \draw[thinmember] (n5) -- (n8);
  \foreach \p in {n1,...,n8}{\node[node dot] at (\p) {};}
  \pic at (n1) {pin}; \pic at (n5) {roller};
  \node[above right] at (n1) {1}; \node[above right] at (n2) {2}; \node[below] at (n3) {3}; \node[below] at (n4) {4}; \node[above left] at (n5) {5};
  \node[above] at (n6) {6}; \node[above] at (n7) {7}; \node[above] at (n8) {8};
  \newcommand{\mem}[2]{\node[circle,draw,fill=white,inner sep=1pt,font=\scriptsize] at #1 {#2};}
  \mem{(3,0)}{1} \mem{(9,0)}{2} \mem{(15,0)}{3} \mem{(21,0)}{4}
  \mem{(3,4)}{5} \mem{(6,4)}{6} \mem{(9,4)}{7} \mem{(12,4)}{8}
  \mem{(13.7,5.7)}{9} \mem{(16.3,5.7)}{10} \mem{(18,4)}{11} \mem{(21,4)}{12}
  \mem{(9,8)}{13} \mem{(15,8)}{14}
  \draw[load,very thick] (n2) -- ++(0,-1.6) node[right] {64\,kN};
  \draw[dashed] (n6) -- (-1.0,8);
  \draw[dim] (-1.0,0) -- (-1.0,8) node[midway,fill=white,inner sep=1pt,rotate=90] {32\,m};
  \draw[dim] (0,-2.3) -- (24,-2.3) node[midway,fill=white,inner sep=1pt] {4@24\,m $=$ 96\,m};
  \draw (0,-2.6) -- (0,-2.0); \draw (24,-2.6) -- (24,-2.0);
\end{tikzpicture}

\end{fig}
\begin{itemize}
  \item $E = 300~\mathrm{kN/cm^2} = 3{,}000{,}000~\mathrm{kN/m^2}$ for members 1, 2, 3, 4, 13, and 14
  \item $E = 200~\mathrm{kN/cm^2} = 2{,}000{,}000~\mathrm{kN/m^2}$ for members 5, 6, 7, 8, 9, 10, 11, 12
  \item Cross section area $= 0.10~\mathrm{m^2}$ for members 1, 2, 3, 4, 13, and 14
  \item Cross section area $= 0.20~\mathrm{m^2}$ for members 6, 8, and 11
  \item Cross section area $= 0.30~\mathrm{m^2}$ for members 5, 7, 9, 10 and 12
\end{itemize}
The contents of \texttt{2Dtruss.ipt}:
\begin{lstlisting}[language={},numbers=none,basicstyle=\small\ttfamily]
2-D truss example
kN
m

* NNOD  NBC  NMAT  NSEC  ITP  NNE  IFORCE
   8    14    2     3     2    2     1

* COOR (m)
1    0.   0.
2   24.   0.
3   48.   0.
4   72.   0.
5   96.   0.
6   24.  32.
7   48.  32.
8   72.  32.

* NFIX (u, v)
1   -1  -1
2    0   0
3    0   0
4    0   0
5   -1  -1
6    0   0
7    0   0
8    0   0

* External Load (kN): Px, Py
1   0.    0.
2   0.  -64.
3   0.    0.
4   0.    0.
5   0.    0.
6   0.    0.
7   0.    0.
8   0.    0.

* IDBC
1   1 2 1 1 0
2   2 3 1 1 0
3   3 4 1 1 0
4   4 5 1 1 0
5   1 6 2 3 0
6   2 6 2 2 0
7   3 6 2 3 0
8   3 7 2 2 0
9   4 7 2 3 0
10  3 8 2 3 0
11  4 8 2 2 0
12  5 8 2 3 0
13  6 7 1 1 0
14  7 8 1 1 0

* PROP
1   3000000. 0. 0. 0. 0.
2   2000000. 0. 0. 0. 0.

* SECT
1   0.10 0. 0. 0. 0.
2   0.20 0. 0. 0. 0.
3   0.30 0. 0. 0. 0.
\end{lstlisting}

\section{Set up global degrees of freedom (equation numbering)}
Functions: \texttt{IDMAT}, \texttt{MEMDOF} and \texttt{SEMIBAND}.
\begin{enumerate}[leftmargin=*]
  \item \texttt{IDMAT}: Equation numbering --- transform \texttt{NFIX(NDN,NNOD)} into \texttt{IDND(NDN,NNOD)} and obtain \texttt{NEQ}.

    Example: 2-D truss (\texttt{NDN}=2)
    \begin{fig}
      % 2-D truss example for IDMAT (6 nodes, 9 members)
\begin{tikzpicture}[scale=0.8,font=\small]
  \coordinate (n1) at (0,0); \coordinate (n2) at (0,3); \coordinate (n3) at (3.5,0); \coordinate (n4) at (3.5,3);
  \coordinate (n5) at (7,0); \coordinate (n6) at (7,3);
  \draw[member] (n1) -- (n2) -- (n4) -- (n3) -- (n1); \draw[member] (n2) -- (n3); \draw[member] (n3) -- (n5) -- (n6) -- (n4); \draw[member] (n3) -- (n6);
  \foreach \p in {n1,...,n6}{\node[node dot] at (\p) {};}
  \pic[rotate=90] at (n2) {pin}; \pic[rotate=90] at (n1) {roller};
  \node[below right] at (n1) {1}; \node[above right] at (n2) {2}; \node[below] at (n3) {3}; \node[above] at (n4) {4}; \node[below] at (n5) {5}; \node[above] at (n6) {6};
  \newcommand{\mem}[2]{\node[circle,draw,fill=white,inner sep=1.5pt,font=\scriptsize] at #1 {#2};}
  \mem{(-0.5,1.5)}{1} \mem{(1.75,-0.5)}{2} \mem{(2.3,1.8)}{3} \mem{(1.75,3.5)}{4}
  \mem{(4.0,1.5)}{5} \mem{(5.25,-0.5)}{6} \mem{(5.8,1.2)}{7} \mem{(5.25,3.5)}{8} \mem{(7.5,1.5)}{9}
\end{tikzpicture}

    \end{fig}
    \begin{center}
      \begin{minipage}{0.47\linewidth}\centering
        \texttt{NFIX(NDN,NNOD)}\\[4pt]
        \begin{tabular}{|c|c|c|c|c|c|c|}\hline
              & 1    & 2    & 3 & 4 & 5 & 6\\\hline
          $u$ & $-1$ & $-1$ & 0 & 0 & 0 & 0\\\hline
          $v$ & 0    & $-1$ & 0 & 0 & 0 & 0\\\hline
        \end{tabular}
      \end{minipage}
      \begin{minipage}{0.47\linewidth}\centering
        \texttt{IDND(NDN,NNOD)}\\[4pt]
        \begin{tabular}{|c|c|c|c|c|c|c|}\hline
              & 1    & 2    & 3 & 4 & 5 & 6\\\hline
          $u$ & $-1$ & $-2$ & 2 & 4 & 6 & 8\\\hline
          $v$ & 1    & $-3$ & 3 & 5 & 7 & 9\\\hline
        \end{tabular}
      \end{minipage}
    \end{center}

  \item \texttt{MEMDOF}: create member d.o.f.\ table, \texttt{LM(NDE,NBC)}.
    \begin{itemize}
      \item Use connectivities \texttt{IDBC(1,IB)} and \texttt{IDBC(2,IB)} (the \texttt{IB}th element) and the nodal d.o.f.\ table, \texttt{IDND(NDN,NNOD)}, to create the member d.o.f.\ table, \texttt{LM(NDE,NBC)}.
    \end{itemize}
    \begin{center}
      \texttt{LM(NDE,NBC)}\\[4pt]
      \begin{tabular}{|c|c|c|c|c|c|c|c|c|c|}\hline
                 & 1    & 2    & 3 & 4 & 5 & 6 & 7 & 8 & 9 (\texttt{NBC})\\\hline
        1        & $-1$ & $-1$ &   &   &   &   & 2 &   & \\\hline
        2        & 1    & 1    &   &   &   &   & 3 &   & \\\hline
        3        & $-2$ & 2    &   &   &   &   & 8 &   & \\\hline
        4 (\texttt{NDE}) & $-3$ & 3 &   &   &   &   & 9 &   & \\\hline
      \end{tabular}
    \end{center}

  \item \texttt{SEMIBAND}: calculate the semi-band width (including diagonal) \texttt{NSBAND}. $[\Kmat]$ can be stored in banded form.
    \begin{center}
      \begin{minipage}{0.5\linewidth}\centering
        $\Kmat$ (\texttt{NEQ},\texttt{NEQ})\\[4pt]
        \renewcommand{\arraystretch}{1.3}
        \begin{tabular}{|c|c|c|c|c|c|}\hline
          $K_{11}$ & $K_{12}$ & $K_{13}$ &          &          &         \\\hline
          $K_{12}$ & $K_{22}$ & $K_{23}$ &          &          &         \\\hline
          $K_{13}$ & $K_{23}$ & $K_{33}$ & $K_{34}$ &          &         \\\hline
                   &          & $K_{34}$ & $K_{44}$ & $K_{45}$ & $K_{46}$\\\hline
                   &          &          & $K_{45}$ & $K_{55}$ & $K_{56}$\\\hline
                   &          &          & $K_{46}$ & $K_{56}$ & $K_{66}$\\\hline
        \end{tabular}
      \end{minipage}
      \begin{minipage}{0.42\linewidth}\centering
        $\widetilde{\Kmat}$ (\texttt{NEQ},\texttt{NSBAND})\\[4pt]
        \renewcommand{\arraystretch}{1.3}
        \begin{tabular}{|c|c|c|}\hline
          $K_{11}$ & $K_{12}$ & $K_{13}$\\\hline
          $K_{22}$ & $K_{23}$ &         \\\hline
          $K_{33}$ & $K_{34}$ &         \\\hline
          $K_{44}$ & $K_{45}$ & $K_{46}$\\\hline
          $K_{55}$ & $K_{56}$ &         \\\hline
          $K_{66}$ &          &         \\\hline
        \end{tabular}
      \end{minipage}
    \end{center}
    $K(i,j)$, $j\geq i$ is stored in $\widetilde{K}(i,\,j-i+1)$; for example,
    \begin{align*}
      K(2,3) &\;\rightarrow\; \widetilde{K}(2,2)\\
      K(4,6) &\;\rightarrow\; \widetilde{K}(4,3)
    \end{align*}
    Calculate \texttt{NSBAND} from the \texttt{LM(NDE,NBC)} matrix columnwise using the algorithm below:
    \begin{equation}
      \texttt{NSBAND} = \max_{i=1,\texttt{NBC}}\Bigl(\max_{j=1,\texttt{NDE}}LM(j,i) - \min_{j=1,\texttt{NDE}}LM(j,i) + 1\Bigr)\quad\text{for } LM(j,i) > 0
    \end{equation}
\end{enumerate}

\section{Assemble external loads into the load vector}
Function \texttt{LOAD}: assemble \texttt{EXLD(NDN,NNOD)} into \texttt{GLOAD(NEQ)}.

\topic{Algorithm:}
\begin{lstlisting}[language=Matlab,numbers=none]
GLOAD = zeros(NEQ,1);
for j = 1:NNOD
    for i = 1:NDN
        if IDND(i,j) > 0
            GLOAD(IDND(i,j)) = GLOAD(IDND(i,j))+EXLD(i,j);
        end
    end
end
\end{lstlisting}

\section{Form the global stiffness matrix and add the equivalent member forces to the load vector}
Functions: \texttt{FORMKP} (\texttt{ROTATION}, \texttt{ELKE}).

For each element, perform the following steps to add its stiffness matrix to the banded global stiffness matrix.
\begin{enumerate}
  \item Calculate element transformation (rotation) matrix $\Tmat_{\texttt{NDE}\times\texttt{NDE}}$. (function \texttt{ROTATION})
  \item Calculate element stiffness matrix referred to the local (elemental) coordinate system, $[k_l]_{\texttt{NDE}\times\texttt{NDE}} = \kmat_l$. (function \texttt{ELKE})
  \item Calculate element stiffness matrix $\kmat$ referred to the global (structural) coordinate system using $\kmat = \Tmat^\tr\kmat_l\Tmat$. ($\kmat_l =$ matrix \texttt{EE}, $\kmat =$ matrix \texttt{ELK})
  \item Assemble $\kmat_{\texttt{NDE}\times\texttt{NDE}}$ into the stiffness matrix $\Kmat_{\texttt{NEQ}\times\texttt{NEQ}}$.
\end{enumerate}
Assembling process: $k_{mn}\;\rightarrow\;K_{IJ}$.

\topic{Algorithm (for the \texttt{IB}th element):}
\begin{lstlisting}[language=Matlab,numbers=none]
% Get element DOFs
LDOF = find(LM(:,IB)>0);
GDOF = LM(LDOF,IB);
% Transform the element stiffness matrix from the local axes
% to global axes : EE --> ELK
ELK = T'*EE*T;
% Assemble the global element stiffness matrix to form the global
% stiffness matrix GLK
GLK(GDOF,GDOF) = GLK(GDOF,GDOF) + ELK(LDOF,LDOF);
\end{lstlisting}

For each element, perform the following three steps to add its equivalent member forces to the load vector.
\begin{enumerate}
  \item Calculate element transformation matrix $\Tmat$. (function \texttt{ROTATION})
  \item Calculate equivalent member forces using $\F_g = -\Tmat^\tr\F_l$, where $\F_l$ (\texttt{FEF(NDE,NBC)}) is the fixed-end forces; $\F_g$ is referred to as equivalent member forces.
    \begin{fig}
      % Fixed-end forces (local) -> equivalent member forces (global)
\begin{tikzpicture}[scale=0.9,font=\footnotesize]
  \def\ang{45}
  % ---- F_l ----
  \begin{scope}
    \coordinate (A) at (0,0); \coordinate (B) at (\ang:3.2);
    \draw[member] (A) -- (B);
    % distributed load
    \foreach \t in {0.1,0.25,...,0.95}{\draw[->] ($(A)!\t!(B)+(\ang+90:0.7)$) -- ($(A)!\t!(B)+(\ang+90:0.1)$);}
    \draw[load] ($(A)+(\ang+90:0.9)$) -- ($(A)+(\ang+90:0.15)$) node[pos=0,left] {10};
    \draw[load] ($(B)+(\ang+90:0.9)$) -- ($(B)+(\ang+90:0.15)$) node[pos=0,right] {$-10$};
    \draw[->,thick] ($(A)+(-0.35,-0.3)$) arc (-120:60:0.35); \node[left] at ($(A)+(-0.3,-0.2)$) {200};
    \draw[->,thick] ($(B)+(0.35,0.3)$) arc (60:240:0.35); \node[right] at ($(B)+(0.25,0.35)$) {$-200$};
    \node at ($(A)+(1.6,-0.7)$) {$\F_l$};
  \end{scope}
  \node at (4.4,2.2) {$\Longrightarrow$};
  % ---- -F_l ----
  \begin{scope}[xshift=5.6cm]
    \coordinate (A) at (0,0); \coordinate (B) at (\ang:3.2);
    \draw[member] (A) -- (B);
    \draw[load] ($(A)+(\ang+90:0.15)$) -- ($(A)+(\ang+90:0.9)$) node[left] {$-10$};
    \draw[load] ($(B)+(\ang+90:0.15)$) -- ($(B)+(\ang+90:0.9)$) node[right] {10};
    \draw[->,thick] ($(A)+(0.35,-0.3)$) arc (-60:-240:0.35); \node[left] at ($(A)+(-0.3,-0.2)$) {$-200$};
    \draw[->,thick] ($(B)+(-0.35,0.3)$) arc (120:-60:0.35); \node[right] at ($(B)+(0.25,0.35)$) {200};
    \node at ($(A)+(1.6,-0.7)$) {$-\F_l$};
  \end{scope}
  \node at (10.2,2.2) {$\Longrightarrow$};
  % ---- -T^T F_l (global components) ----
  \begin{scope}[xshift=11.6cm]
    \coordinate (A) at (0,0); \coordinate (B) at (\ang:3.2);
    \draw[member] (A) -- (B);
    \draw[load] (A) -- ++(0,0.9) node[left] {$-10\frac{1}{\sqrt2}$};
    \draw[load] (A) -- ++(0.9,0) node[right] {$10\frac{1}{\sqrt2}$};
    \draw[load] (B) -- ++(0,0.9) node[left] {$-10\frac{1}{\sqrt2}$};
    \draw[load] (B) -- ++(0.9,0) node[below right] {$10\frac{1}{\sqrt2}$};
    \draw[->,thick] ($(A)+(0.35,-0.3)$) arc (-60:-240:0.35); \node[left] at ($(A)+(-0.35,-0.35)$) {$-200$};
    \draw[->,thick] ($(B)+(-0.35,0.3)$) arc (120:-60:0.35); \node[right] at ($(B)+(0.3,0.45)$) {200};
    \node at ($(A)+(1.6,-1.1)$) {$-\Tmat^\tr\F_l$};
  \end{scope}
\end{tikzpicture}

    \end{fig}
  \item Assemble $\F_g$ into \texttt{GLOAD(NEQ)} using the algorithm below (for the \texttt{IB}th element).
\end{enumerate}
\begin{lstlisting}[language=Matlab,numbers=none]
Steps 2 and 3:
EFEQ = -T'*FEF(:,IB);
GLOAD(GDOF) = GLOAD(GDOF) + EFEQ(LDOF);
\end{lstlisting}

\section{Solve for the nodal unknowns}
(Chapter~\ref{ch:solve}; function \texttt{SOLVE})

\section{Calculate member end forces (force recovery)}
(function \texttt{FORCE})

For each element (member), the member end forces can be calculated as follows:
\begin{enumerate}
  \item Use \texttt{LM(NDE,IB)} (the \texttt{IB}th element) and the global displacement vector \texttt{DISP(NEQ)} to find the element displacement vector in the global coordinate system, $\dvec_g$.
  \item Call function \texttt{ROTATION} to obtain the element transformation matrix $\Tmat$.
  \item Obtain the element displacement vector in the local coordinate system, $\dvec_l$ (dimension: $\texttt{NDE}\times\texttt{NDE}$) using $\dvec_l = \Tmat\dvec_g$.
  \item Call function \texttt{ELKE} to obtain the element stiffness matrix $\kmat_l$.
  \item Calculate the element forces (member end forces) using $\mat{f}_l = \kmat_l\dvec_l + \F_l$.
\end{enumerate}
